% ================================================================== Chapitre 10
\chapter{Formulaire et données}
\label{ch:formulaire}

\section*{Formulaire}
\addcontentsline{toc}{section}{Formulaire}

\renewcommand{\arraystretch}{1.55}
\begin{xltabular}{\linewidth}{>{\raggedright\arraybackslash}p{4.3cm} X}
\toprule
\multicolumn{2}{l}{\sffamily\bfseries\color{navy} Potentiel chimique et enthalpie libre} \\ \midrule
Enthalpie libre & $G = H - TS = U + PV - TS$ \\
Identité (composition fixée) & $\dd G = V\,\dd P - S\,\dd T$ \\
Mélange & $\dd G = V\,\dd P - S\,\dd T + \sum_i\mu_i\,\dd n_i$, \quad $G = \sum_in_i\mu_i$ \\
Potentiel chimique & $\mu_i = \left(\partial G/\partial n_i\right)_{T,P,n_{j\neq i}} = \mu_i^\circ(T) + RT\ln a_i$ \\
Activités & gaz : $P_i/P^\circ$ ; soluté : $c_i/c^\circ$ ; solvant, solide ou liquide pur : $1$ \\ \midrule
\multicolumn{2}{l}{\sffamily\bfseries\color{navy} Grandeurs de réaction} \\ \midrule
Grandeur de réaction & $\Delta_rX = (\partial X/\partial\xi)_{T,P} = \sum_i\nu_iX_{m,i}$ \quad ; \quad $\Delta_rG = \sum_i\nu_i\mu_i$ \\
Grandeur standard & $\Delta_rX^\circ(T) = \sum_i\nu_iX^\circ_{m,i}$ \quad ; \quad $\Delta_rG^\circ = \Delta_rH^\circ - T\Delta_rS^\circ$ \\
Loi de Hess & $\Delta_rH^\circ = \sum_i\nu_i\Delta_fH^\circ_i$, \quad $\Delta_rG^\circ = \sum_i\nu_i\Delta_fG^\circ_i$, \quad $\Delta_rS^\circ = \sum_i\nu_iS^\circ_{m,i}$ \\
ESR & $\Delta_fX^\circ(\text{corps simple dans son ESR}) = 0$ \quad (mais $S^\circ_m \neq 0$) \\
Combinaison linéaire & $(3) = a(1) + b(2) \Rightarrow \Delta_rX^\circ_3 = a\Delta_rX^\circ_1 + b\Delta_rX^\circ_2$, \quad $K^\circ_3 = (K^\circ_1)^a(K^\circ_2)^b$ \\
Ellingham & $\Delta_rH^\circ$, $\Delta_rS^\circ$ constantes hors changement d'état ; $T_i = \Delta_rH^\circ/\Delta_rS^\circ$ \\ \midrule
\multicolumn{2}{l}{\sffamily\bfseries\color{navy} Premier principe} \\ \midrule
Monobare isotherme & $Q = \Delta H = \xi_F\,\Delta_rH^\circ$ \\
Monobare adiabatique & $\xi_F\,\Delta_rH^\circ(T_I) + \big(\sum_jn_jC^\circ_{p,m,j}\big)(T_F - T_I) = 0$ \\ \midrule
\multicolumn{2}{l}{\sffamily\bfseries\color{navy} Évolution et équilibre ($T$, $P$ fixées)} \\ \midrule
Création d'entropie & $T\,\delta S_c = -\Delta_rG\,\dd\xi \geq 0$ \quad ; \quad $\dd G = -T\,\delta S_c \leq 0$ \\
Enthalpie libre de réaction & $\Delta_rG = \Delta_rG^\circ + RT\ln Q_r = RT\ln(Q_r/K^\circ)$, \quad $Q_r = \prod_ia_i^{\nu_i}$ \\
Constante d'équilibre & $K^\circ = \exp(-\Delta_rG^\circ/RT)$ \quad ; \quad équilibre : $Q_r = K^\circ$, $\Delta_rG = 0$ \\
Sens d'évolution & $\Delta_rG<0 \Leftrightarrow Q_r<K^\circ$ : sens direct \\
Van~'t~Hoff & $\dfrac{\dd\ln K^\circ}{\dd T} = \dfrac{\Delta_rH^\circ}{RT^2}$ \quad ; \quad $\ln\dfrac{K^\circ_2}{K^\circ_1} = -\dfrac{\Delta_rH^\circ}{R}\left(\dfrac1{T_2} - \dfrac1{T_1}\right)$ \\
Gaz : effet de $P$ & $Q_r = \prod n_i^{\nu_i}\big(P/(n_{\text{tot,gaz}}P^\circ)\big)^{\Delta\nu_{\text{gaz}}}$ \\ \midrule
\multicolumn{2}{l}{\sffamily\bfseries\color{navy} Oxydoréduction et piles} \\ \midrule
Nernst & $E = E^\circ + \dfrac{RT}{\gamma F}\ln\dfrac{a_{\mathrm{Ox}}^\alpha}{a_{\mathrm{Red}}^\beta} \approx E^\circ + \dfrac{0{,}06}{\gamma}\log\dfrac{a_{\mathrm{Ox}}^\alpha}{a_{\mathrm{Red}}^\beta}$ \\
Réaction rédox & $\Delta_rG = -\gamma F(E_1 - E_2)$, \quad $\Delta_rG^\circ = -\gamma F(E^\circ_1 - E^\circ_2)$ \\
Constante & $\log K^\circ = \gamma(E^\circ_1 - E^\circ_2)/0{,}06$ \quad à $\SI{298}{K}$ \\
Demi-équation & $\Delta_{1/2}G^\circ = -\gamma FE^\circ$ \quad (on additionne les $\Delta G^\circ$, jamais les $E^\circ$) \\
Travail électrique & $\Delta G \leq W_{\text{élec}}$ \quad ; \quad $W_{\max,\text{fourni}} = -\Delta G$ \\
Tension à vide & $e = E_\oplus - E_\ominus = -\Delta_rG/(\gamma F) = e^\circ - \dfrac{0{,}06}{\gamma}\log Q_r$ \quad ; \quad pile usée : $e = 0$ \\
Capacité & $q_{\max} = \gamma F\xi_{\max}$, \quad $\SI{1}{A.h} = \SI{3600}{C}$, \quad $W\approx q_{\max}\,e$ \\
\bottomrule
\end{xltabular}

\begin{pieges}
\begin{enumerate}[leftmargin=1.8em]
  \item Oublier les entropies molaires des corps simples dans $\Delta_rS^\circ$.
  \item Mélanger kJ et J dans $\Delta_rG^\circ = \Delta_rH^\circ - T\Delta_rS^\circ$.
  \item Utiliser le signe de $\Delta_rG^\circ$ (au lieu de $\Delta_rG$) pour prévoir le sens d'évolution.
  \item Oublier les espèces spectatrices (\ce{N2} de l'air, eau, calorimètre) dans le calcul d'une température finale.
  \item Additionner des potentiels standard au lieu des enthalpies libres $-\gamma FE^\circ$.
  \item Écrire $\Delta_rG^\circ = +\gamma FE^\circ$ : le signe moins est indispensable (une pile, spontanée, a $e^\circ > 0$ et $\Delta_rG^\circ<0$).
  \item Écrire $G = U - PV - TS$ : c'est $G = U + PV - TS$.
  \item Croire qu'ajouter un gaz inerte n'a jamais d'effet : c'est vrai à $V$ fixé, pas à $P$ fixée.
\end{enumerate}
\end{pieges}

\section*{Données thermodynamiques à 298 K}
\addcontentsline{toc}{section}{Données thermodynamiques}

\begin{table}[H]
\centering
\renewcommand{\arraystretch}{1.18}
\small
\begin{tabular}{lrrrr}
\toprule
Constituant & $\Delta_fH^\circ$ ($\mathrm{kJ\cdot mol^{-1}}$) & $S^\circ_m$ ($\mathrm{J\cdot K^{-1}\cdot mol^{-1}}$) & $\Delta_fG^\circ$ ($\mathrm{kJ\cdot mol^{-1}}$) & $C^\circ_{p,m}$ ($\mathrm{J\cdot K^{-1}\cdot mol^{-1}}$) \\ \midrule
\ce{N2(g)}      & 0       & 191,6 & 0       & 29,1 \\
\ce{H2(g)}      & 0       & 130,7 & 0       & 28,8 \\
\ce{O2(g)}      & 0       & 205,2 & 0       & 29,4 \\
\ce{Cl2(g)}     & 0       & 223,1 & 0       & 33,9 \\
C (graphite)    & 0       & 5,7   & 0       & 8,5 \\
C (diamant)     & 1,9     & 2,4   & 2,9     & 6,1 \\
\ce{NH3(g)}     & $-45{,}9$   & 192,8 & $-16{,}4$  & 35,1 \\
\ce{CH4(g)}     & $-74{,}6$   & 186,3 & $-50{,}5$  & 35,7 \\
\ce{CH3Cl(g)}   & $-81{,}9$   & 234,6 & --      & -- \\
\ce{HCl(g)}     & $-92{,}3$   & 186,9 & $-95{,}3$  & 29,1 \\
\ce{CO(g)}      & $-110{,}5$  & 197,7 & $-137{,}2$ & 29,1 \\
\ce{CO2(g)}     & $-393{,}5$  & 213,8 & $-394{,}4$ & 37,1 \\
\ce{H2O(l)}     & $-285{,}8$  & 70,0  & $-237{,}1$ & 75,3 \\
\ce{H2O(g)}     & $-241{,}8$  & 188,8 & $-228{,}6$ & 33,6 \\
\ce{CaCO3(s)}   & $-1207{,}6$ & 91,7  & $-1129{,}1$ & 83,5 \\
\ce{CaO(s)}     & $-634{,}9$  & 38,1  & $-603{,}3$ & 42,0 \\
\ce{ZnO(s)}     & $-350{,}5$  & 43,7  & $-320{,}5$ & -- \\
\ce{Ag2O(s)}    & $-31{,}1$   & 121,3 & $-11{,}2$  & -- \\
\ce{HO-(aq)}    & $-230{,}0$  & --    & --      & -- \\
\bottomrule
\end{tabular}
\caption{Données thermodynamiques à $\SI{298}{K}$ (valeurs usuelles des tables). Convention pour les ions : $\Delta_fH^\circ(\ce{H+(aq)}) = 0$.}
\label{tab:donnees}
\end{table}

\begin{table}[H]
\centering
\renewcommand{\arraystretch}{1.18}
\small
\begin{tabular}{lr@{\hspace{2.5em}}lr}
\toprule
Couple & $E^\circ$ (V) & Couple & $E^\circ$ (V) \\ \midrule
\ce{O2/H2O}       & 1,23  & \ce{Cu^{2+}/Cu+}  & 0,16 \\
\ce{Hg2^{2+}/Hg}  & 0,797 & \ce{H+/H2}        & 0 (référence) \\
\ce{Fe^{3+}/Fe^{2+}} & 0,77 & \ce{Fe^{2+}/Fe}  & $-0{,}44$ \\
\ce{Cu^{2+}/Cu}   & 0,34  & \ce{Zn^{2+}/Zn}   & $-0{,}76$ \\
\bottomrule
\end{tabular}
\caption{Potentiels standard à $\SI{298}{K}$ utilisés dans ce cours.}
\label{tab:potentiels}
\end{table}
